\uuid{v2V0}
\exo7id{4299}
\titre{exo7 4299}
\auteur{quercia}
\organisation{exo7}
\datecreate{2010-03-12}
\isIndication{false}
\isCorrection{false}
\chapitre{Calcul d'intégrales}
\sousChapitre{Intégrale impropre}
\cpge{mpsi-mp2i.raisonnement-vocabulaire-ensembliste.a, pcsi.raisonnement-vocabulaire-ensembliste.a, ptsi.raisonnement-vocabulaire-ensembliste.a, mpsi-mp2i.integration.c, pcsi.integration.b, ptsi.integration.b, bcpst1.integration-segment.a, bcpst1.integration-segment.b, mp-mpi.integration-intervalle.b, bcpst2.integrales-generalisees}
\module{Analyse}
\niveau{L1}
\difficulte{}

\contenu{
\texte{
Soit $f : {[0,+\infty[} \to {\R}$ uniformément continue telle que
$ \int_{t=0}^{+\infty} f(t)\,d t$ converge.
}
\begin{enumerate}
    \item \question{Montrer que $f(t) \to 0$ lorsque $t\to{+\infty}$(raisonner par l'absurde).}
    \item \question{Si $f$ est positive, montrer que $ \int_{t=0}^{+\infty} f^2(t)\,d t$ converge.}
    \item \question{Donner un contre-exemple si $f$ n'est pas de signe constant.}
\end{enumerate}
}
